% Clase 4 --- La suma en la recta: construcción por paralelogramos (§8.1).
% O, I, Y, X marcados sobre la recta principal R (coordenadas: O=0, I=1,
% Y=1.8, X=3). El punto suma Z=X+Y=4.8 se lee sobre R, pero se justifica
% con un paralelogramo auxiliar levantado a altura h=1.4: se traslada el
% vector u=OY (mismo vector, dos copias) para que arranque en X, y su
% extremo, bajado verticalmente, da Z sobre la recta.
\begin{center}
\begin{tikzpicture}[scale=1]

  % ---- recta principal R ----
  \draw[figaxis] (-0.6,0) -- (5.5,0) node[right, figlbl] {$R$};

  % ---- puntos sobre R ----
  \node[figptocerrado] (O) at (0,0)   {};
  \node[figptocerrado] (I) at (1,0)   {};
  \node[figptocerrado] (Y) at (1.8,0) {};
  \node[figptocerrado] (X) at (3,0)   {};
  \node[figptocerrado] (Z) at (4.8,0) {};

  \node[figlbl, below=3pt] at (O) {$O$};
  \node[figlbl, below=3pt] at (I) {$I$};
  \node[figlbl, below=3pt] at (Y) {$Y$};
  \node[figlbl, below=3pt] at (X) {$X$};
  \node[figlbl, below=3pt] at (Z) {$X{+}Y$};

  % ---- paralelogramo auxiliar (levantado a altura 1.4) ----
  \node[figptoabierto] (Yp) at (1.8,1.4) {};
  \node[figptoabierto] (Zp) at (4.8,1.4) {};

  % lados que trasladan el vector u = OY hasta arrancar en X
  \draw[figvec] (O) -- (Yp) node[figlblaux, midway, left=1pt] {$u$};
  \draw[figvec] (X) -- (Zp) node[figlblaux, midway, right=1pt] {$u$};

  % lado superior: copia trasladada del segmento OX (paralela y de igual
  % longitud), cierra el paralelogramo O, X, Z', Y'
  \draw[figline] (Yp) -- (Zp);

  % bajada auxiliar: se proyecta Z' verticalmente sobre R para leer X+Y
  \draw[figaux] (Zp) -- (Z);

\end{tikzpicture}\\[0.5em]
{\small\itshape Construcción de $x+y$ por traslación de vectores: el
paralelogramo $O,X,(X{+}Y)',Y'$ (levantado por claridad; en la
construcción real vive sobre la propia recta $R$, donde el
``paralelogramo'' degenera en un segmento) traslada el vector
$u=\overrightarrow{OY}$ hasta que arranca en $X$; su extremo, proyectado de
vuelta sobre $R$, es el punto $X{+}Y$. La suma es una operación
\textbf{lineal}: no requiere conocer la unidad $I$ (marcada aquí solo como
referencia).}
\end{center}
